\section*{v2.1.0 Structural Audit and Re-Expansion Core}
\addcontentsline{toc}{section}{v2.1.0 Structural Audit and Re-Expansion Core}
\begin{quote}
\textbf{Normalization note.} This block is retained as a historical audit precursor. Its main value is to preserve re-expansion discipline and proof-memory traceability; the normalized document no longer treats it as a separate generative core, but as supporting governance for the later unfolding phases.
\end{quote}

This release adds a further consolidation layer to the companion monolith. The
new emphasis is not merely on carrying closure information forward, but on
making explicit how later re-expansion remains auditable after several stages of
compression. Earlier versions already introduced closure ladders, interface
filters, ledger surfaces, and proof-memory discipline. Version~v2.1.0 now
adds an \emph{audit core}: every admissible compressed statement should retain a
recoverable expansion route that can be checked against the same closure class
from which it descended.

\subsection*{1. Auditability as structural maturity}
A compressed statement becomes structurally mature only when it is possible to
state not just where it came from, but how it may be unfolded again without
changing its admissible class. This requirement sharpens the earlier ledger
ideas. The ledger tracked which closure path licensed a statement; the audit
core now asks whether that path can be re-opened in a controlled way.

\paragraph{Audit Proposition A (licensed re-expansion).}
A normalized claim is fully admissible only if there exists a bounded
re-expansion route from the compressed form back to a compatible lemma, tool,
probe, or closure ancestry class. A statement that cannot be re-opened under
controlled transport may still be heuristically useful, but it has not yet
reached final structural normalization.

\paragraph{Proof sketch.}
Closure ancestry alone prevents arbitrary insertion of unsupported projection
claims, but it does not yet guarantee that later readers can recover the path.
The present strengthening adds a re-opening condition. The compressed surface
must not merely summarize the proof memory; it must preserve enough structure
that the same memory can be reconstructed under admissible reversal.

\subsection*{2. The audit ladder}
The working audit ladder of this release is
\[
\begin{aligned}
&\text{local triolic constraint} \Longrightarrow \text{HC compensation} \Longrightarrow \text{closure confinement}\\
&\Longrightarrow \text{probe/interface filtration} \Longrightarrow \text{lemma or table compression}\\
&\Longrightarrow \text{readout statement} \Longrightarrow \text{controlled re-expansion}.
\end{aligned}
\]
The final arrow is new in emphasis. It makes explicit that a readout is not the
terminal endpoint of the proof architecture. Instead, it is a temporarily stable
surface from which the architecture should remain recoverable.

\paragraph{Audit Proposition B (endpoints remain reversible).}
A readout that has lost all reversible relation to its closure ladder should be
classified as provisional. The stronger the reversible link to its generating
path, the stronger its status inside the monolith.

\subsection*{3. Probe-audit discipline}
The probe channel receives an additional task in v2.1.0. It is no longer
only a selection or filtration device; it also serves as an audit checkpoint.
Probe-selected modes should leave a downstream trace that remains visible in the
compressed lemma/table/readout layers. This makes the probe sector a bridge
between admissible dynamics and later recoverable proof memory.

\paragraph{Audit Proposition C (probe traces survive compression).}
If a configuration is admissibly selected in the orthogonal probe channel, then
its later compressed representations should preserve a recognizable probe trace.
If no such trace remains, downstream formulations must be marked as reduced,
heuristic, or not yet fully synchronized.

\subsection*{4. Appendix as an audit surface}
This version also sharpens the role of the appendix. Appendix entries, tables,
and named tools now act as \emph{audit surfaces}. Each item should indicate not
only what it abbreviates, but how it can be re-opened into the same structural
class. The appendix therefore becomes more than a memory store or ledger sheet;
it becomes a reversible interface between compressed exposition and detailed
proof ancestry.

\paragraph{Audit Proposition D (appendix entries carry reopening metadata).}
A mature appendix entry should point both backward and forward: backward into
its closure ancestry and forward into an admissible re-expansion channel. This
dual role prevents tables from drifting into detached summaries.

\subsection*{5. Consequence for the next mathematical thickening}
With v2.1.0 the document takes another step toward a larger mathematically
expanded companion. Future lemma growth should now be judged not only by local
correctness, but also by whether it improves auditability across compression
levels. New material is strongest when it reduces the gap between compressed
statement, appendix row, probe trace, and recoverable closure ancestry.

\paragraph{Working rule for subsequent versions.}
Whenever a new mathematical block is added, it should be possible to specify
which closure class licenses it, which compressed interfaces abbreviate it, and
which controlled re-expansion route reconstructs it. Material satisfying all
three demands strengthens the proof architecture in the intended structure-first
sense.



\clearpage
